\chapter{\IfLanguageName{dutch}{Stand van zaken}{State of the art}}%
\label{ch:stand-van-zaken}

% Tip: Begin elk hoofdstuk met een paragraaf inleiding die beschrijft hoe
% dit hoofdstuk past binnen het geheel van de bachelorproef. Geef in het
% bijzonder aan wat de link is met het vorige en volgende hoofdstuk.

% Pas na deze inleidende paragraaf komt de eerste sectiehoofding.
Om inzicht te verkrijgen in deze complexe interactiestromen is een analysemethode nodig die processen kan reconstrueren op basis van operationele data. Deze sectie bespreekt daarom de fundamentele concepten van process mining, de rol van event logs en de specifieke uitdaging die ontstaan bij de toepassing van process mining binnen omnichannelcontactcenters.

De evolutie van traditionele callcenters naar omnichannelcontactcenters heeft geleid tot een fundamentele verandering in hoe klanterinteracties worden verwerkt en geanalyseerd \autocite{kalra2024}. Waar vroegere contactcenters voornamelijk voice-gebaseerd waren, integreren moderne platformen zoals Genesys gemiddeld meerdere communicatiekanalen, waaronder telefonie, webchat, social-messagingplatformen en e-mail. Deze kanaaldiversiteit creëert niet-lineaire en sterk variabele klantpaden, wat de complexiteit van interactiestromen aanzienlijk verhoogt. \textcite{kalra2024} toont aan dat organisaties hierdoor steeds vaker geconfronteerd worden met gefragmenteerde operationele logs die niet ontworpen zijn voor end-to-end flowanalyse. Dit bemoeilijkt het identificeren van inefficiënties, foutpaden en drop-offs en benadrukt de nood aan methoden die volledige klantpaden kunnen reconstrueren en analyseren \autocite{kalra2024}.

Om deze complexe en niet-lineaire klantpaden te kunnen analyseren is een methode nodig die ruwe logs omzet naar interpreteerbare procesmodellen. Process mining vormt hiervoor de standaardbenadering. In de volgende sectie worden de belangrijkste concepten en technieken binnen process mining besproken.

Process mining is een datagedreven analysetechniek die procesmodellen afleidt uit event logs en zo inzicht biedt in de werkelijke uitvoering van processen \autocite{vanDerAalst2016}. Het vormt een brug tussen data science en business process management en maakt het mogelijk om processen te ontdekken, te vergelijken en te verbeteren.
Process mining omvat traditioneel drie hoofdactiviteiten: process discovery, conformance checking en process enhancement \autocite{vanDerAalst2016,Zelst2020}. Bij process discovery wordt op basis van event logs automatisch een procesmodel opgebouwd zonder voorafgaande kennis van het proces. Conformance checking vergelijkt vervolgens het ontdekte proces met een bestaand referentiemodel om afwijkingen, uitzonderingen of niet-conforme uitvoeringen te identificeren. Process enhancement  is het proces om met de verkregen inzichten verdere optimalisatie, bijvoorbeeld door bottlenecks, wachttijden of inefficiënte processtappen zichtbaar te maken \autocite{vanDerAalst2016}.
De kwaliteit van process-mininganalyse hangt sterk af van de beschikbare event logs. Een event log bevat typisch informatie over een specifieke case, de uitgevoerde activiteit, een timestap en bijkomende attributen zoals gebruiker, kanaal of systeem \autocite{vanDerAalst2016}. In traditionele bedrijfsprocessen zijn deze logs vaak relatief gestructureerd en afkomstig uit één centraal systeem. Binnen omnichannelcontactcenters daarentegen worden interacties verspreid geregistreerd over verschillende platformen en applicaties waardoor verspreide datasets ontstaan.

Deze uiteenlopende datasets brengen verschillende uitdagingen met zich mee voor process mining \autocite{Zelst2020}. Ten eerste verschillende communicatiekanalen die verschillen in de manier waarop interacties worden geregistreerd. Telefoniesystemen genereren doorgaans andere events, terwijl chat- en messagingplatformen eerder conversatiegebaseerde logs produceren. Daarnaast verlopen interacties binnen omnichannelomgevingen vaak asynchroon en parallel over meerdere communicatiekanalen tegelijk. Hierdoor wordt het reconstrueren van consistente end-to-end klantpaden aanzienlijk complexer \autocite{kalra2024}.

Een bijkomende uitdaging is de correlatie van gebeurtenissen afkomstig uit verschillende systemen naar één uniforme procesinstantie. Klantinteracties kunnen verspreid zijn over CRM-systemen, ticketingplatformen, chatdiensten en telefonie-infrastructurenen, waarbij niet altijd een gemeenschappelijke identificatie beschikbaar is. Hierdoor moeten technieken toegepast worden om die gebeurtenissen te koppelen op basis van tijdsvensters, metadata of contextuele informatie \autocite{Zelst2020}. Onvoldoende eventcorrelatie kan leiden tot onvolledige of vertekende procesmodellen waardoor analyses minder betrouwbaar worden.
Om deze uitdaging te verwerken werden binnen process mining verschillende technieken ontwikkeld voor het ontdekken van procesmodellen uit event logs. Eén van de meest gekende benaderingen is process discovery waarbij algoritmen automatisch processtructruren reconstrueren op basis van geregistreerde gebeurtenissen. Bekende discovery-algoritmen zijn onder andere de Alpha Miner, Heuristic Miner en Inductive Miner \autocite{vanDerAalst2016}. Vooral de inductive Minder wordt frequent toegepast omdat deze robuuster omgaat met ruis, parallelisme en complexe procesvariaties die typisch voorkomen binnen servicegerichte omgevingen zoals omnichannelcontactcenters.

Binnen omnichannelcontactcenters biedt process mining organisaties de mogelijkheid om volledige klanttrajecten inzichtelijk te maken en operationele inefficiënties te identificeren. Door interacties over verschillende communicatiekanalen heen te analyseren kunnen bottlenecks, herhaalde contactpogingen en problematische overdrachtsmomenten systematisch opgespoord worden. Hierdoor vormt process mining een waardevolle techniek voor het optimaliseren van klantinteracties en het verbeteren van de algemene customer experience \autocite{kalra2024}.

Ondanks deze mogelijkheden blijft de praktische implementatie binnen omnichannelomgevingen uitdagend door de hoge variabiliteit van interactiestromen en de nood aan consistente eventcorrelatie tussen systemen. Verdere onderzoek en aangepaste analysetechnieken blijven daarom noodzakelijk om process mining effectief toe te passen binnen moderne contactcenterarchitecturen.


% Dit hoofdstuk bevat je literatuurstudie. De inhoud gaat verder op de inleiding, maar zal het onderwerp van de bachelorproef *diepgaand* uitspitten. De bedoeling is dat de lezer na lezing van dit hoofdstuk helemaal op de hoogte is van de huidige stand van zaken (state-of-the-art) in het onderzoeksdomein. Iemand die niet vertrouwd is met het onderwerp, weet nu voldoende om de rest van het verhaal te kunnen volgen, zonder dat die er nog andere informatie moet over opzoeken \autocite{Pollefliet2011}.

% Je verwijst bij elke bewering die je doet, vakterm die je introduceert, enz.\ naar je bronnen. In \LaTeX{} kan dat met het commando \texttt{$\backslash${textcite\{\}}} of \texttt{$\backslash${autocite\{\}}}. Als argument van het commando geef je de ``sleutel'' van een ``record'' in een bibliografische databank in het Bib\LaTeX{}-formaat (een tekstbestand). Als je expliciet naar de auteur verwijst in de zin (narratieve referentie), gebruik je \texttt{$\backslash${}textcite\{\}}. Soms is de auteursnaam niet expliciet een onderdeel van de zin, dan gebruik je \texttt{$\backslash${}autocite\{\}} (referentie tussen haakjes). Dit gebruik je bv.~bij een citaat, of om in het bijschrift van een overgenomen afbeelding, broncode, tabel, enz. te verwijzen naar de bron. In de volgende paragraaf een voorbeeld van elk.

% \textcite{Knuth1998} schreef een van de standaardwerken over sorteer- en zoekalgoritmen. Experten zijn het erover eens dat cloud computing een interessante opportuniteit vormen, zowel voor gebruikers als voor dienstverleners op vlak van informatietechnologie~\autocite{Creeger2009}.

% Let er ook op: het \texttt{cite}-commando voor de punt, dus binnen de zin. Je verwijst meteen naar een bron in de eerste zin die erop gebaseerd is, dus niet pas op het einde van een paragraaf.

% \begin{figure}
%   \centering
%   \includegraphics[width=0.8\textwidth]{grail.jpg}
%   \caption[Voorbeeld figuur.]{\label{fig:grail}Voorbeeld van invoegen van een figuur. Zorg altijd voor een uitgebreid bijschrift dat de figuur volledig beschrijft zonder in de tekst te moeten gaan zoeken. Vergeet ook je bronvermelding niet!}
% \end{figure}

% \begin{listing}
%   \begin{minted}{python}
%     import pandas as pd
%     import seaborn as sns

%     penguins = sns.load_dataset('penguins')
%     sns.relplot(data=penguins, x="flipper_length_mm", y="bill_length_mm", hue="species")
%   \end{minted}
%   \caption[Voorbeeld codefragment]{Voorbeeld van het invoegen van een codefragment.}
% \end{listing}

% \lipsum[7-20]

% \begin{table}
%   \centering
%   \begin{tabular}{lcr}
%     \toprule
%     \textbf{Kolom 1} & \textbf{Kolom 2} & \textbf{Kolom 3} \\
%     $\alpha$         & $\beta$          & $\gamma$         \\
%     \midrule
%     A                & 10.230           & a                \\
%     B                & 45.678           & b                \\
%     C                & 99.987           & c                \\
%     \bottomrule
%   \end{tabular}
%   \caption[Voorbeeld tabel]{\label{tab:example}Voorbeeld van een tabel.}
% \end{table}

